\appendix
\section{Disease parameters}
\label{app:sources}

Every parameter value is taken from the literature, not tuned. Each clinical rate is the
reciprocal of a measured duration of the exemplar disease ($\mathrm{rate}=1/\mathrm{duration}$);
each basic reproduction number $R_0$ comes from a separate epidemiological study; and the
transmission rate $\beta$ is derived from both, $\beta=R_0\gamma$ (or $\beta=R_0(\gamma+\kappa)$
for the isolation type). Table~\ref{tab:params-values} lists the value of every rate per type, and
Table~\ref{tab:sources} gives the source for each. Notably, measles and Ebola have similar
infectious periods ($\sim$8--10~d) yet $R_0$ differs roughly sevenfold ($\approx$15 vs
$\approx$2): the infectious period does not set $R_0$; together with $\beta$ it produces it.

\begin{table*}[t]
\centering
\footnotesize
\caption{Parameter values per disease type (rates per day; $\mathrm{rate}=1/\mathrm{duration}$,
with the duration in parentheses). No $\langle k^2\rangle/\langle k\rangle$ correction is applied,
as nodes are well-mixed cities rather than individuals. The endemic type's $\gamma$ and $R_0$ are
illustrative: with no immunity and a long, variable untreated duration, the regime matters, not
the exact number. For the lethal type $\gamma_Q\approx\gamma\approx0.10$.}
\label{tab:params-values}
\begin{tabular}{@{}ll cccccc@{}}
\toprule
\# & Type $\cdot$ model & $\gamma$ & $\sigma$ & $\kappa$ & $\omega$ & $\mu$ & $R_0$ \\
\midrule
1 & Immunizing $\cdot$ SIR    & 0.13 ($\sim$8\,d)   & ---                 & 0                 & 0                          & $\approx$0  & 12--18 \\
2 & Latent $\cdot$ SEIR       & 0.13 ($\sim$7--8\,d) & 0.20 ($\sim$5.1\,d) & 0                 & 0                          & $\sim$0.01  & 2--3.3 \\
3 & Lethal $\cdot$ SEIQRD     & 0.10 ($\sim$10\,d)  & 0.09 ($\sim$11.4\,d) & 0.20 ($\sim$5\,d) & 0                          & 0.71        & 1.5--2.5 \\
4 & Endemic $\cdot$ SIS       & $\sim$0.01--0.05    & ---                 & 0                 & $\infty$                   & 0           & low ($\gtrsim$1) \\
5 & Recurrent $\cdot$ SEIRS   & 0.21 ($\sim$4.8\,d) & 0.70 ($\sim$1.4\,d) & 0                 & $\sim$0.001 (2--7\,yr)     & $\sim$0.001 & 1.3--1.8 \\
\bottomrule
\end{tabular}
\end{table*}

\begin{table*}[t]
\centering
\footnotesize
\caption{Source for every disease-type parameter. $\gamma$, $\sigma$, $\kappa$ come from a measured
clinical duration of the exemplar; $R_0$ from a separate epidemiological study; $\beta$ is derived
from both.}
\label{tab:sources}
\begin{tabular}{@{}llll>{\raggedright\arraybackslash}p{7.4cm}@{}}
\toprule
Disease (model) & Value & Measured quantity & Source (paraphrased) \\
\midrule
Measles (SIR) & $\gamma=0.13$ & infectious $\sim$8~d & contagious 4~d before to 4~d after rash~\cite{cdc:measles} \\
              & $R_0=12$--$18$ & $R_0$ & ``often cited to be 12--18'' (review of 58 estimates)~\cite{guerra:measles} \\
              & $\omega=0$ & immunity & natural infection confers life-long immunity~\cite{cdc:measles} \\
\addlinespace
COVID (SEIR) & $\sigma=0.20$ & incubation 5.1~d & ``mean incubation period 5.1 days (95\% CI 4.5--5.8)''~\cite{lauer:incubation} \\
             & $\gamma=0.13$ & infectious $\sim$1~wk & infectiousness peaks near onset, low by $\sim$7~d~\cite{he:shedding} \\
             & $R_0=2$--$3.3$ & $R_0$ & ``pooled $R_0$ 3.32 (95\% CI 2.81--3.82)''~\cite{alimohamadi:covidr0} \\
\addlinespace
Ebola (SEIQRD) & $\sigma=0.09$ & incubation 11.4~d & ``incubation period was 11.4 days''~\cite{whoebola2014} \\
               & $\gamma,\kappa$ struct. & onset$\rightarrow$hospital $\sim$5~d & SEIHFR; community / hospital / burial transmission~\cite{legrand:ebola} \\
               & $\mu=0.71$ & case-fatality & ``case fatality rate 70.8\% (95\% CI 69--73)''~\cite{whoebola2014} \\
               & $R_0=1.5$--$2.5$ & $R_0$ (SEIR fit) & 1.51 Guinea, 2.53 Sierra Leone, 1.59 Liberia~\cite{althaus:ebolar0} \\
\addlinespace
Gonorrhea (SIS) & SIS, $\omega=\infty$ & immunity, latency & ``negligible protective immunity, negligible latent period''~\cite{hethcote:gonorrhea} \\
                & $\gamma\approx0.03$ & untreated duration & median untreated pharyngeal gonorrhea $\approx$16 weeks~\cite{barbee:gonorrhea} \\
\addlinespace
Influenza (SEIRS) & $\sigma=0.70$ & incubation 1.4~d & ``median incubation period for influenza A 1.4 days''~\cite{lessler:incubation} \\
                  & $\gamma=0.21$ & shedding 4.8~d & ``viral shedding averaged 4.80 days (95\% CI 4.31--5.29)''~\cite{carrat:influenza} \\
                  & $\omega\approx0.001$ & immunity & influenza A immunity wanes in 2--7~yr via antigenic drift~\cite{woolthuis:influenza} \\
                  & $R_0=1.3$--$1.8$ & $R_0$ (review) & median $R\approx1.80$ for the 1918 pandemic (IQR 1.47--2.27)~\cite{biggerstaff:influenzar0} \\
\bottomrule
\end{tabular}
\end{table*}

\section{Implementation architecture}
\label{app:arch}

The interactive simulator and the programmatic staged-protocol runs share one engine,
configuration system, and results pipeline (Figure~\ref{fig:arch}), so the demonstration and the
study cannot drift apart. (This figure is a working schematic and will be refined.)

\begin{figure}[t]
  \centering
  \includegraphics[width=\columnwidth]{F5-architecture.pdf}
  \caption{System architecture: one shared engine behind the interactive simulator and the
  programmatic runs, writing to a common results pipeline.}
  \label{fig:arch}
\end{figure}
